\section{The finite-set bound for the weighted norms}
\label{sw-section}

We can now combine the arithmetic cycles, finite-field
functionals, homology bound, and detection of global
powers.  This proves assertion~\ref{sc-few}, and also
its version with the smaller scalar exponent anticipated
in (SC19).  The classical rational homology theorems used
in Section~\ref{sb-section} and the qualitative
Chebotarev theorem used in Section~\ref{sd-section}
remain the explicit classical inputs to this argument.

Let \(F/L\) be a cyclic extension of number fields with
generator \(\gamma\).  Suppose that \(L\) is totally
imaginary and that
\[
 \nu\ge c\ge1,\qquad c\ge2\text{ if }\ell=2,
 \qquad [F:L]=\ell^{\nu-c},
 \qquad \gamma\zeta=\zeta^{1+\ell^c},
 \tag{SW1}
\]
where \(\ell\) is prime and \(\zeta\in F\) is a
primitive \(q\)-th root of unity, with \(q=\ell^\nu\).
Put
\[
 \kappa=1+\ell^c,\quad E\ge1,\quad
 m=2E+1<\ell^c,\quad a=\kappa^E,\quad A=\mathbb Z/q.
\]
Call \(\alpha\in F^\times\) admissible if there is a
fractional divisor \(\mathfrak d\) on \(F\) such that
\[
             v_P(\alpha)=\mathfrak d(\gamma P)
                                      -a\mathfrak d(P)
                  \quad\text{for all finite primes }P.
\]
As before, define
\[
 \mathcal N(\alpha)=\prod_{j=0}^{[F:L]-1}
                              (\gamma^j\alpha)^{a^j},
 \quad
 c_E=\prod_{\substack{0\le j\le2E\\j\ne E}}
                              (\kappa^E-\kappa^j),
 \quad
 b_*=2E(c+1)+2v_\ell(E!).
 \tag{SW2}
\]

\begin{theorem}\label{sw-few}
Let \(b\ge b_*\) be an integer.  If \(T\subset F^\times\)
is a finite set of admissible elements and
\[
 \left(\frac{\mathcal N(\alpha)}{\mathcal N(\alpha')}
              \right)^{\ell^b}\in(F^\times)^q
          \quad\Longrightarrow\quad\alpha=\alpha'
                       \qquad(\alpha,\alpha'\in T),
 \tag{SW3}
\]
then
\[
 \#T\le q^{r_2(L)}
           t_L(4E+6,2E+1)t_L(4E+6,2E).
 \tag{SW4}
\]
This conclusion does not require \(2E+1<r_2(L)\).
\end{theorem}
\begin{proof}
For each \(\alpha\in T\), choose a divisor witnessing
admissibility.  Theorem~\ref{sa-class} supplies a size
\(N_\alpha\), a cycle
\[
 y_\alpha\in C_m(\operatorname{GL}_{N_\alpha}
                           (\mathcal O_L);A),
\]
and a nonzero \(b_\alpha\in\mathcal O_L\) with its
simultaneous reduction property.  The cycle and
\(b_\alpha\) are chosen before any residue field or
finite-field functional.  Set
\[
                  b_0=\ell\prod_{\alpha\in T}b_\alpha
                         \in\mathcal O_L\setminus\{0\}.
\]
Let \(\mathcal P\) be the primes \(P\) of \(F\)
such that \(P\nmid b_0\) and the natural map
\(\mathcal O_L\to k_P:=\mathcal O_F/P\) is
surjective.

For \(P\in\mathcal P\), reduction is injective on
\(\mu_q\), by Lemma~\ref{sd-roots-reduce}, since
\(P\nmid\ell\).  Thus the reduction \(\overline\zeta\)
has order \(q\), and \(q\mid\#k_P^\times\).
Choose a generator \(g_P\) of the cyclic group
\(k_P^\times\), and define
\[
      v_P:k_P^\times\longrightarrow A,
                       \qquad v_P(g_P^j)=j\pmod q.
\]
The divisibility just proved makes this well-defined.
It is surjective and
\[
                         \ker v_P=(k_P^\times)^q.
 \tag{SW5}
\]
Indeed an exponent maps to zero precisely when it is
divisible by \(q\).

Theorem~\ref{sf-functional} constructs an additive
family \(\theta_P\) with the scalar-prism value
specified by \(v_P\) and \(\overline\zeta\).
Reduce matrices over \(\mathcal O_L\) to \(k_P\)
and apply this family.  Corollary~\ref{sf-primitive-family}
gives linear maps
\[
 \Phi_{P,N}:C_m(\operatorname{GL}_N(\mathcal O_L);A)
                                     \longrightarrow A
\]
which vanish on boundaries, are stable on cycles,
and multiply their value by \(d\) on \(d\) repeated
direct sums.  The defining property of the chosen
arithmetic cycle is
\[
 \Phi_{P,N_\alpha}(y_\alpha)
     =c_E\,v_P\bigl(\overline{\mathcal N(\alpha)}\bigr)
                              \qquad(P\in\mathcal P).
 \tag{SW6}
\]
The residue on the right exists and is nonzero:
\(\mathcal N(\alpha)\) is a unit in
\(\mathcal O_F[1/b_\alpha]\) by
Theorem~\ref{sa-class}.

Take the \(A\)-module
\(B=\prod_{P\in\mathcal P}A\) and let
\(\Phi_N=(\Phi_{P,N})_{P\in\mathcal P}\).
All three properties of the family hold
coordinate by coordinate.  No finiteness assumption
on this product is needed in Theorem~\ref{sb-upper}.
Define
\[
 f:T\longrightarrow B,
      \qquad f(\alpha)=\ell^{2E}\Phi_{N_\alpha}(y_\alpha).
 \tag{SW7}
\]
That theorem bounds \(\#f(T)\) by the right-hand side
of (SW4), since \(m-1=2E\) and \(2m+4=4E+6\).
We prove that \(f\) is injective.

Suppose \(f(\alpha)=f(\alpha')\), and put
\(\beta=\mathcal N(\alpha)/\mathcal N(\alpha')\).
For every \(P\in\mathcal P\), (SW6) gives
\[
                    \ell^{2E}c_E v_P(\overline\beta)=0
                                                \quad\text{in }A.
\]
Lemma~\ref{sc-vandermonde} gives
\(v_\ell(c_E)=2Ec+2v_\ell(E!)\), and proves the
exact cancellation of the factor prime to \(\ell\).
Consequently
\[
 v_P(\overline\beta^{\,\ell^b})
       =\ell^b v_P(\overline\beta)=0
                      \qquad(P\in\mathcal P),
\]
because \(b\ge b_*\).
Equation (SW5) says that the residue of
\(\beta^{\ell^b}\) is a \(q\)-th power at every
such prime.  The element is a unit there, as noted
after (SW6).

Corollary~\ref{sd-source}, applied to
\(\beta^{\ell^b}\) and the nonzero integral element
\(b_0\), now gives
\(\beta^{\ell^b}\in(F^\times)^q\).
The separating hypothesis (SW3) implies
\(\alpha=\alpha'\).  Thus \(f\) is injective, and
the bound for \(\#f(T)\) proves (SW4).
\end{proof}

\begin{corollary}\label{sw-source-few}
Assertion~\ref{sc-few}, with its original exponent
\[
                 b=2E+(2E+1)(c+2E+1),
\]
holds.  Its strengthened version with \(b=b_*\)
also holds under the same hypotheses.
\end{corollary}
\begin{proof}
The valuation of \(E!\) is
\(\sum_{j\ge1}\lfloor E/\ell^j\rfloor\le E\),
by counting the multiples of each prime power and
bounding the resulting geometric series.
Hence
\[
 b_*\le 2E+2Ec+2E
          \le 2E+(2E+1)(c+2E+1).
\]
Apply Theorem~\ref{sw-few} with each of the two
exponents.  The source's additional hypothesis
\(2E+1<r_2(L)\) is permitted but unnecessary here.
\end{proof}

\begin{corollary}\label{sw-coarse-consequence}
With the classical reciprocity and duality results
stated earlier, the bounds (SP12) and (SP13) hold
without assuming assertion~\ref{sc-few} separately.
More explicitly, if \(L\) is totally imaginary,
\(E\ge1\), and \(\ell\) is prime, with
\(\mu_\ell\subset L\) for odd \(\ell\) and
\(\mu_4\subset L\) for \(\ell=2\), then
\[
 \begin{split}
 \log\#\mathcal S_\ell(L,E)
   &\le \exp\bigl((2E[L:\mathbb Q])^{86}
                          \log(2|\operatorname{disc}L|)\bigr),\\
 \log\#H^2(G_{L,S},\mathbb Z_\ell(E+1))
   &\le \exp\bigl((2E[L:\mathbb Q])^{88}
                          \log(2|\operatorname{disc}L|)\bigr),
 \end{split}
\]
where \(S\) is the set of primes above \(\ell\).
Both groups are finite.
\end{corollary}
\begin{proof}
Corollary~\ref{sw-source-few} supplies the remaining
hypothesis of Corollary~\ref{sp-conditional-coarse}.
Its proof combines the already proved homology torsion
estimate, class-number bound, unit comparison,
cyclotomic descent, and strict Selmer comparison.
No additional quantitative comparison is introduced
at this last step.
\end{proof}

The independent algebraic \(K\)-theory route of
Section~\ref{kt-section} gives the stronger coarse
exponent \(14\).  The purpose of the present argument
is different: it completes the conventional proof of
the source's actual finite-set comparison, while
improving its scalar exponent.  This does not by
itself prove the remaining theta and period-map
arguments or the final analytic assertion.


